Two preliminaries. First, $\mathfrak{sl}(2,\C)$ is simple: with the standard basis $e, h, f$, an ideal containing $x = ae + bh + cf \neq 0$ contains $[e,[e,x]] = -2c\,e$ and $[f,[f,x]] = -2a\,f$, so it contains $e$ or $f$ unless $a = c = 0$, in which case it contains $h$ and then $[h,e] = 2e$; and from any one of $e, f, h$ the brackets $[e,f] = h$, $[h,e] = 2e$, $[h,f] = -2f$ produce the other two. Second, the ideals of $\mathfrak{l} = \mathfrak{g}^+\oplus\mathfrak{g}^-$ are $0$, $\mathfrak{g}^+$, $\mathfrak{g}^-$, $\mathfrak{l}$: if an ideal $I$ is not contained in $\mathfrak{g}^-$, its projection to $\mathfrak{g}^+$ is a nonzero ideal of $\mathfrak{g}^+$, hence all of $\mathfrak{g}^+$, and then $\mathfrak{g}^+ = [\mathfrak{g}^+, \mathfrak{g}^+] = [\mathfrak{g}^+, I] \subseteq I$ (the $\mathfrak{g}^-$-components of $I$ commute with $\mathfrak{g}^+$); symmetrically for $\mathfrak{g}^-$.

(1) The Levi factor stabilises the rays of $\nO$ and $\nInf$, so by equivariance of $\Psi$ it preserves the planes $\Spi$ and $\Som$, and $\mathfrak{l}$ acts on both. The map $\omega\otimes\pi \mapsto \omega\wedge\pi$ is equivariant for the induced action on $\Lambda^2\T$ (the action for which $\Psi$ is equivariant) and surjective onto the four-dimensional mixed slot $\Som\wedge\Spi = \Psi(W_\C)$ (Proposition \claimref{correspondence.skew-lines-splittings-hyperbolic-planes-one}), so by dimension $\Som\otimes\Spi \cong W_\C$ as $\mathfrak{l}$-modules. The action of $\mathfrak{l}$ on $W_\C$ is faithful: the bivector $e_i\wedge e_j$ of an orthonormal pair acts as $\pm$ the elementary antisymmetric matrix, and the six are linearly independent. Let $\mathfrak{k}_\omega, \mathfrak{k}_\pi \subseteq \mathfrak{l}$ be the kernels of the actions on $\Som$, $\Spi$: ideals. Each is nonzero, since $\dim\mathfrak{l} = 6 > 4 = \dim\mathfrak{gl}(2,\C)$. Neither is $\mathfrak{l}$: if $\mathfrak{l}$ annihilated $\Spi$ it would act on $W_\C \cong \Som\otimes\Spi$ through $\Som$ alone, with image of dimension at most $4 < 6$, contradicting faithfulness; likewise for $\Som$. So $\mathfrak{k}_\omega, \mathfrak{k}_\pi \in \{\mathfrak{g}^+, \mathfrak{g}^-\}$, and they are different, since a common kernel would annihilate $\Som\otimes\Spi = W_\C$. Say $\mathfrak{k}_\pi = \mathfrak{g}^a$ and $\mathfrak{k}_\omega = \mathfrak{g}^b$ with $\{a, b\} = \{+,-\}$. The kernel of $\mathfrak{g}^b$ acting on $\Spi$ is an ideal of $\mathfrak{g}^b$, hence $0$ or $\mathfrak{g}^b$, and not $\mathfrak{g}^b$ (then $\mathfrak{l} = \mathfrak{g}^a + \mathfrak{g}^b$ would annihilate $\Spi$); so $\mathfrak{g}^b$ acts faithfully on $\Spi$, and likewise $\mathfrak{g}^a$ on $\Som$. The verification script exhibits the assignment in the Fock model.

(2) Keep the labelling of (1). A faithful action of $\mathfrak{sl}(2,\C)$ on a $2$-plane $S$ is an isomorphism onto $\mathfrak{sl}(S)$: the image is three-dimensional and consists of traceless operators, because $\mathfrak{sl}(2,\C) = [\mathfrak{sl}(2,\C), \mathfrak{sl}(2,\C)]$ and commutators are traceless, and $\dim\mathfrak{sl}(S) = 3$. An $\mathfrak{l}$-equivariant map $f: W_\C\otimes\Spi \cong \Som\otimes\Spi\otimes\Spi \to \Som$ is therefore a linear map that commutes with $\mathfrak{sl}(\Som)$ acting on the first factor (through $\mathfrak{g}^a$, which acts trivially on the $\Spi$ factors) and is annihilated by $\mathfrak{sl}(\Spi)$ acting as a derivation on the last two factors (through $\mathfrak{g}^b$, trivial on $\Som$). Regard $f$ as a bilinear map $\Spi\times\Spi \to \End(\Som)$, $(\pi,\pi') \mapsto (\omega \mapsto f(\omega\otimes\pi\otimes\pi'))$. Its values commute with $\mathfrak{sl}(\Som)$, hence are scalars (an operator commuting with $E_{12}$, $E_{21}$ and $E_{11} - E_{22}$ is scalar), so $f$ is a bilinear form $\beta$ on $\Spi$ with $\beta(X\pi, \pi') + \beta(\pi, X\pi') = 0$ for all $X \in \mathfrak{sl}(\Spi)$. In a basis $\pi_1, \pi_2$ with $h = \operatorname{diag}(1,-1)$ and $e = E_{12}$: $h$-invariance gives $2\beta(\pi_1,\pi_1) = 0 = -2\beta(\pi_2,\pi_2)$, and $e$-invariance applied to $(\pi_2, \pi_2)$ gives $\beta(\pi_1,\pi_2) + \beta(\pi_2,\pi_1) = 0$. So $\beta$ is a multiple of the determinant form, which is invariant; the space is one-dimensional. The other case is the same with the roles of the planes exchanged.
